\chapter{Just Enough Complex Analysis}

Chapter~1 treated $z$ as a formal symbol, manipulated according to algebraic rules
with no thought of convergence or magnitude.
The real power of analytic combinatorics comes from promoting $z$ to a genuine
complex variable and extracting the asymptotic growth of the coefficients
$[z^n] A(z)$ from the analytic structure of the function $A$.
This chapter develops the minimum complex analysis required to make that move
rigorous and computational.
The reader who has seen complex numbers but not complex functions will find
everything needed here; the reader who has taken a full complex analysis course
may skim and use this chapter as a reference for notation.
The standard treatment of all topics below is \cite{FlajoletSedgewick2009}, Part~A.

\section{Complex Numbers and Convergence}

Recall that $\mathbb{C} = \{x + iy : x, y \in \mathbb{R}\}$ with the convention
$i^2 = -1$.
The \emph{modulus} of $z = x + iy$ is $|z| = \sqrt{x^2 + y^2}$, and the
\emph{argument} $\arg z$ is the angle $\theta \in (-\pi, \pi]$ such that
$x = |z|\cos\theta$, $y = |z|\sin\theta$.
Euler's formula
\[
  e^{i\theta} = \cos\theta + i\sin\theta
\]
allows every nonzero complex number to be written in \emph{polar form}
$z = r e^{i\theta}$ with $r = |z| > 0$ and $\theta = \arg z$.
Multiplication becomes transparent in polar form: $r_1 e^{i\theta_1} \cdot
r_2 e^{i\theta_2} = r_1 r_2 \, e^{i(\theta_1+\theta_2)}$, so moduli multiply
and arguments add.

Convergence in $\mathbb{C}$ is measured by the modulus: a sequence $(z_k)$
converges to $L$ if $|z_k - L| \to 0$.
A series $\sum_k w_k$ \emph{converges absolutely} if $\sum_k |w_k| < \infty$,
and absolute convergence implies convergence.
The geometric series identity
\[
  \sum_{k=0}^{N} w^k = \frac{1 - w^{N+1}}{1 - w}, \quad w \neq 1,
\]
shows that $\sum_{k=0}^{N} w^k \to 1/(1-w)$ as $N \to \infty$ whenever
$|w| < 1$, since $|w|^{N+1} \to 0$.
We will use this repeatedly.

\section{Power Series and Radius of Convergence}

Let $A(z) = \sum_{n \ge 0} a_n z^n$ be a formal power series with coefficients
$a_n \in \mathbb{C}$.
Treating $z$ as a complex variable, the set of $z$ for which this series converges
absolutely is a disk.

\begin{theorem}[Cauchy--Hadamard]
  The power series $A(z) = \sum_{n \ge 0} a_n z^n$ converges absolutely for every
  $z$ with $|z| < R$ and diverges for every $z$ with $|z| > R$, where the
  \emph{radius of convergence} $R$ is given by
  \[
    \frac{1}{R} = \limsup_{n \to \infty} |a_n|^{1/n}.
  \]
\end{theorem}

\begin{proof}[Proof sketch]
Fix $z$ with $|z| < R$. Then $|z| < 1/\limsup |a_n|^{1/n}$, so there exists
$\delta > 0$ with $|z| \le (1-\delta)/\limsup |a_n|^{1/n}$. By the definition
of $\limsup$, for all sufficiently large $n$ we have
$|a_n|^{1/n} \le (1/R) + \epsilon$ for any $\epsilon > 0$; choosing $\epsilon$
small enough that $(1/R+\epsilon)|z| < 1$, we get
$|a_n z^n| = (|a_n|^{1/n}|z|)^n \le r^n$ for some $r < 1$ and all large $n$.
The series $\sum |a_n z^n|$ is then bounded above by a convergent geometric
series, so $\sum a_n z^n$ converges absolutely.

Conversely, if $|z| > R$, then $\limsup |a_n|^{1/n} |z| > 1$, so
$|a_n z^n| > 1$ for infinitely many $n$, and the terms do not tend to zero.
The series diverges.
\end{proof}

The combinatorial import of Cauchy--Hadamard is immediate: the radius of
convergence captures the \emph{exponential growth rate} of $(a_n)$.
If $R > 0$ is finite, then $a_n$ grows (roughly) like $R^{-n}$, in the sense that
$\limsup |a_n|^{1/n} = R^{-1}$ (and for the well-behaved sequences arising from
combinatorics, the ordinary limit exists and equals $R^{-1}$).
More precisely, $a_n = R^{-n} \cdot \theta(n)$ where $\theta(n)$ is
\emph{subexponential}: $\limsup |\theta(n)|^{1/n} = 1$.
Determining the fine structure of $\theta(n)$ is precisely what asymptotic
analysis does.

\begin{example}
  For $B(z) = 1/(1-2z) = \sum_{n \ge 0} 2^n z^n$, we have $b_n = 2^n$ and
  $|b_n|^{1/n} = 2$ for all $n$, so $R = 1/2$.
  The singularity of $B$ at $z = 1/2$ is directly visible as a pole,
  and the exponential growth rate $2^n = (1/2)^{-n}$ agrees with $R^{-n}$.

  For the Fibonacci generating function $F(z) = z/(1 - z - z^2)$,
  the poles are at the roots of $1 - z - z^2 = 0$, namely $z = (-1 \pm \sqrt{5})/2$.
  The root of smallest modulus is $1/\varphi$ where $\varphi = (1+\sqrt{5})/2
  \approx 1.618$ is the golden ratio, so $R = 1/\varphi$ and
  $f_n \sim C \varphi^n$ for some explicit constant $C$.
\end{example}

\section{Analytic Functions}

A function $f: U \to \mathbb{C}$ defined on an open set $U \subseteq \mathbb{C}$
is \emph{holomorphic} on $U$ (also called \emph{analytic}) if it is
complex-differentiable at every point of $U$.
A fundamental and non-trivial theorem of complex analysis states that the
following conditions are equivalent for a function on an open set:
(1) $f$ is complex-differentiable (the limit $\lim_{h \to 0}(f(z+h)-f(z))/h$
exists in $\mathbb{C}$);
(2) $f$ is infinitely differentiable and satisfies the Cauchy--Riemann equations
$\partial_x u = \partial_y v$, $\partial_y u = -\partial_x v$ where $f = u + iv$;
(3) $f$ is locally representable by a convergent power series around every point.
We will take this equivalence as given.

The intuition behind this equivalence is worth pausing on, even without a proof.
Real differentiability requires only that a function have a well-defined slope when
you approach a point along the real line---from the left or from the right.
Complex differentiability demands something far more stringent: the limit
$(f(z+h)-f(z))/h$ must yield the \emph{same} value no matter what direction $h$
approaches zero in the complex plane---horizontally, vertically, diagonally,
or along any curve whatsoever.
This is an overdetermined system: two partial derivatives of the real and imaginary
parts must simultaneously satisfy the Cauchy--Riemann equations, which is already
a strong constraint, and having this hold at every point is stronger still.
The result is that a function with even one complex derivative everywhere is
automatically infinitely smooth, and its values in any open disk are completely
determined by a convergent power series---a rigidity that has no real analogue.
Real-differentiable functions can have kinks and corners in their higher derivatives;
complex-differentiable functions cannot, because the complex constraint leaves no
room for local irregularity.

The key example: if $A(z) = \sum_{n \ge 0} a_n z^n$ has radius of convergence
$R > 0$, then $A$ is holomorphic on the open disk $|z| < R$.
This follows because power series can be differentiated term by term within
their disk of convergence, and the resulting series again converges.
Conversely, a holomorphic function on $|z| < R$ is completely determined by
its power series at the origin.
The formal generating functions of Chapter~1, as soon as their coefficients
grow at most geometrically, are genuine holomorphic functions on a nontrivial disk.

\begin{definition}
A \emph{singularity} of a function $f$ is a point $z_0 \in \mathbb{C}$ at
which $f$ fails to be holomorphic: there is no convergent power series
representation of $f$ in any neighborhood of $z_0$.
A \emph{dominant singularity} is a singularity of smallest modulus $|z_0|$.
\end{definition}

\section{The Geometric Series as Prototype}

The function $f(z) = 1/(1-z)$ is the model example.
For $|z| < 1$ the series $\sum_{n \ge 0} z^n$ converges absolutely, and
multiplying by $(1-z)$ collapses the partial sums to $1 - z^{N+1} \to 1$, giving
\[
  \frac{1}{1-z} = \sum_{n \ge 0} z^n, \quad |z| < 1.
\]
The function $1/(1-z)$ is holomorphic on all of $\mathbb{C} \setminus \{1\}$:
the formula $1/(1-z)$ makes perfect sense and defines a smooth function at,
say, $z = 2$ or $z = -3$, even though the series $\sum z^n$ diverges there.
In the language of complex analysis, $1/(1-z)$ is the \emph{analytic continuation}
of the series beyond its disk of convergence --- the unique holomorphic function
that agrees with the series where the series converges, but exists on a larger domain.
The series itself, however, is unable to see past the boundary $|z|=1$.
The point $z = 1$ is a \emph{simple pole}: $f(z) = 1/(1-z)$ blows up like
$(1-z)^{-1}$ near $z = 1$, which is the simplest possible singularity.

The lesson is structural: the singularity at $z=1$ is what \emph{prevents}
the radius of convergence from exceeding $1$, because the series cannot
represent a function that blows up.
In general, the radius of convergence of a power series at a point $z_0$
equals the distance from $z_0$ to the nearest singularity of the function
that the series represents.
Every analytic-combinatorics computation ultimately reduces to locating
the nearest singularity of a generating function.

\section{Cauchy's Integral Formula for Coefficients}

The coefficients of a power series can be recovered by a contour integral.
Let $f$ be holomorphic on $|z| < R$ and let $0 < r < R$.
Then
\[
  a_n = [z^n] f(z) = \frac{1}{2\pi i} \oint_{|z| = r} \frac{f(z)}{z^{n+1}}\, dz,
\]
where the contour integral runs counterclockwise around the circle of radius $r$.
This follows from Cauchy's integral formula for derivatives: for a holomorphic
$f$ on an open disk,
\[
  f^{(n)}(0) = \frac{n!}{2\pi i} \oint_{|z|=r} \frac{f(z)}{z^{n+1}}\, dz,
\]
combined with Taylor's theorem $f(z) = \sum_n a_n z^n$ implying
$f^{(n)}(0) = n!\, a_n$.
Dividing both sides by $n!$ gives the coefficient formula.

The crucial observation for asymptotics is that the contour $|z| = r$ can be
\emph{deformed outward}, provided it remains in the region where $f$ is
holomorphic. Suppose $f$ is holomorphic on a disk of radius $r$ (with all
singularities outside $|z| = r$). Write $M(r) = \max_{|z|=r}|f(z)|$. Then
the standard \emph{ML bound} on contour integrals gives
\[
  |a_n| = \frac{1}{2\pi}\biggl|\oint_{|z|=r}\frac{f(z)}{z^{n+1}}\,dz\biggr|
  \le \frac{1}{2\pi}\cdot 2\pi r \cdot \frac{M(r)}{r^{n+1}}
  = \frac{M(r)}{r^n}.
\]
(The factor $2\pi r$ is the length of the circle; $|f(z)/z^{n+1}| \le M(r)/r^{n+1}$
on the contour.) This is an \emph{exponential} upper bound on $|a_n|$: the decay
rate $r^{-n}$ improves as $r$ increases, while $M(r)$ is a constant (depending
on $r$ but not on $n$). The contour can be pushed outward as far as possible
without crossing a singularity of $f$; the singularity forms an impassable
barrier. Taking $r$ just below the distance to the nearest singularity $\rho$
gives $|a_n| \le M(\rho-\epsilon)/(\rho-\epsilon)^n$ for any $\epsilon > 0$,
which is essentially $\rho^{-n}$ times a constant. The dominant contribution to
$a_n$ therefore comes from the behavior of $f$ near its nearest singularity.

This is the geometric heart of singularity analysis: singularities are not
obstacles to computation but rather the \emph{source} of asymptotic information.

\section{Classifying Singularities}

The type of a singularity determines the subexponential correction to the
exponential growth rate.
For a generating function with non-negative real coefficients --- which is the
case for every counting problem --- a classical result (the Pringsheim--Vivanti
theorem) guarantees that the point $z = R$ on the positive real axis is always
a singularity. We write $\rho = R > 0$ for this dominant singularity.
(Generating functions with complex or alternating-sign coefficients can have
their dominant singularity elsewhere in $\mathbb{C}$, but this case rarely
arises in combinatorics.)
The three types that appear most frequently are the following.

\textbf{Poles.}
Consider $f(z) = 1/(1 - z/\rho)^k$ for a positive integer $k$. We build up the
coefficient formula from first principles. For $k=1$, this is the geometric
series: $1/(1-w) = \sum_{n \ge 0} w^n$ with $w = z/\rho$, giving
$[z^n]f = \rho^{-n}$. For $k=2$, differentiate the geometric series:
$\frac{d}{dw}\frac{1}{1-w} = \frac{1}{(1-w)^2} = \sum_{n \ge 1} n\, w^{n-1}
= \sum_{n \ge 0} (n+1) w^n$, so $[z^n]\frac{1}{(1-z/\rho)^2} = (n+1)\rho^{-n}$.
Differentiating again: $\frac{1}{(1-w)^3} = \frac{1}{2}\sum_{n \ge 0}(n+2)(n+1)w^n$,
so $[z^n] = \binom{n+2}{2}\rho^{-n}$. In general, differentiating $k-1$ times:
\[
  \frac{1}{(1 - z/\rho)^k} = \sum_{n \ge 0} \binom{n + k - 1}{k - 1} \rho^{-n} z^n.
\]
(This can also be derived from the generalized binomial theorem of Chapter~1
with $\alpha = -k$.)

The coefficient $\binom{n+k-1}{k-1}$ is a polynomial of degree $k-1$ in $n$.
To see this, write it out:
\[
  \binom{n+k-1}{k-1} = \frac{(n+k-1)(n+k-2)\cdots(n+1)}{(k-1)!}.
\]
The numerator is a product of $k-1$ factors, each of the form $n + (\text{constant})$,
so it is a polynomial of degree $k-1$ in $n$ with leading term $n^{k-1}$.
Dividing by the constant $(k-1)!$:
\[
  \binom{n+k-1}{k-1} \;\sim\; \frac{n^{k-1}}{(k-1)!} \quad\text{as } n \to \infty.
\]
A simple pole ($k=1$) contributes a pure exponential $\rho^{-n}$; a double
pole ($k=2$) contributes $(n+1)\rho^{-n} \sim n\rho^{-n}$; a triple pole
($k=3$) contributes $\binom{n+2}{2}\rho^{-n} \sim \frac{n^2}{2}\rho^{-n}$;
and so on. Most rational generating functions have poles, and they arise
naturally whenever a linear recurrence has a unique dominant root.

\textbf{Square-root branch points.}
If $f(z) = \sqrt{1 - z/\rho} = (1 - z/\rho)^{1/2}$, then the generalized
binomial theorem gives $f(z) = \sum_{n \ge 0} \binom{1/2}{n} (-z/\rho)^n$.
In Chapter~1 we derived $\binom{1/2}{n} = (-1)^{n-1}/(2^{2n-1}n)\binom{2n-2}{n-1}$,
so $[z^n]f = \binom{1/2}{n}(-1)^n\rho^{-n} = -\rho^{-n}/(2^{2n-1}n)\binom{2n-2}{n-1}$.

It remains to estimate $\binom{2n-2}{n-1}$ for large $n$. This is a central
binomial coefficient (up to a shift), and its asymptotics follow from
Stirling's approximation $m! \approx \sqrt{2\pi m}\,(m/e)^m$. Write $m = n-1$:
\[
  \binom{2m}{m} = \frac{(2m)!}{(m!)^2}
  \approx \frac{\sqrt{2\pi(2m)}\,(2m/e)^{2m}}{(\sqrt{2\pi m}\,(m/e)^m)^2}
  = \frac{\sqrt{4\pi m}\, \cdot\, 2^{2m} m^{2m}/e^{2m}}{2\pi m \,\cdot\, m^{2m}/e^{2m}}
  = \frac{2^{2m}}{\sqrt{\pi m}}.
\]
So $\binom{2n-2}{n-1} \sim 4^{n-1}/\sqrt{\pi(n-1)} \sim 4^{n-1}/\sqrt{\pi n}$
for large $n$. Substituting and noting $4^{n-1}/2^{2n-1} = 2^{2n-2}/2^{2n-1} = 1/2$:
\[
  [z^n]\sqrt{1 - z/\rho} = -\frac{1}{2\sqrt{\pi}}\, n^{-3/2}\, \rho^{-n}\bigl(1 + O(1/n)\bigr).
\]
The exponent $-3/2$ in the polynomial correction $n^{-3/2}$ is the
\emph{fingerprint} of a square-root singularity.
This type arises pervasively in combinatorics because the kernel method and
quadratic equations produce square roots; trees, lattice paths, and
many self-referential structures lead to generating functions with
$\sqrt{1 - z/\rho}$ as a factor.
We will return to this in Chapter~5.

\textbf{Logarithmic singularities.}
If $f(z) = \log(1/(1 - z/\rho))$, then $f(z) = \sum_{n \ge 1} \rho^{-n} z^n / n$,
so $[z^n] f = \rho^{-n}/n$.
The correction factor is $1/n$: it decays to zero (unlike the constant
correction at a simple pole) but more slowly than the $n^{-3/2}$ of a
square-root singularity.
Logarithms appear, for instance, in the generating functions for labelled
structures involving cycles, and in the analysis of certain algorithms.
They are less common than poles or branch points but arise naturally whenever
the symbolic construction includes $\SET$ or $\CYC$ operators.

\textbf{Essential singularities.}
Functions such as $e^{1/(1-z)}$ have an essential singularity at $z = 1$:
they blow up in a highly irregular, non-power-law fashion that cannot be
captured by any finite Laurent expansion.
The elementary transfer-theorem machinery does not apply at essential
singularities, and special methods are required.
Such singularities are rare in the combinatorial generating functions
considered in this book, and we will not develop the necessary tools here.

\section{The Big Picture}

The moral of this chapter is simple to state and powerful to apply.
If $A(z)$ is a generating function whose nearest singularity to the origin
lies at $|z| = \rho$, then $a_n$ grows exponentially like $\rho^{-n}$.
That much is determined purely by the \emph{location} of the singularity.
The \emph{type} of the singularity---pole of order $k$, square-root branch
point, logarithm---controls the subexponential factor $n^{\alpha}$ for some
$\alpha$ determined by the local singular behavior.
Flajolet and Odlyzko's transfer theorem \cite{FlajoletOdlyzko1990}, developed in Chapter~4, makes
this dictionary precise and gives explicit asymptotic expansions to any
desired order of precision \cite{FlajoletSedgewick2009}.

All the complex analysis we use in this book reduces to three steps:
(i) find a useful closed form for the generating function $A(z)$, using the
recurrence method of Chapter~1 or the symbolic method of Chapter~3; (ii) locate the singularities on the boundary
of the disk of convergence; (iii) read the asymptotics off their type using
the classification above.
The deeper machinery of complex analysis---contour deformation through branch
cuts, analytic continuation to Riemann surfaces, the residue theorem in full
generality, the saddle-point method---belongs to a dedicated complex-analysis
course and will appear in this book only when no simpler approach suffices.
The examples of Chapter~3 will make the three-step program concrete before
we prove the transfer theorem in Chapter~4.

\begin{exercise}
Find the radius of convergence of $\sum_{n \ge 1} z^n / n^2$. What is the dominant singularity? (Hint: use Cauchy--Hadamard. Note that $|1/n^2|^{1/n} \to 1$.)
\end{exercise}

\begin{exercise}
Show that $1/(1-z)^3 = \sum_{n \ge 0} \binom{n+2}{2} z^n$ by applying the generalized binomial theorem (Chapter~1) with $\alpha = -3$ and $w = -z$. The singularity at $z = 1$ is a pole of order $3$; verify that $\binom{n+2}{2} = (n+2)(n+1)/2 \sim n^2/2$, consistent with the catalogue entry $n^{k-1}$ for $k = 3$.
\end{exercise}

\begin{exercise}
The generating function $\sum_{n \ge 0} n!\, z^n$ has radius of convergence $R = 0$ (verify this using Cauchy--Hadamard). What does this say about the location of singularities? Can you make sense of the ``function'' this series represents? (This is an example where the formal power series of Chapter~1 is a valid algebraic object but has no analytic avatar.)
\end{exercise}

\section*{Further Reading}

\begin{itemize}
  \item \textbf{P. Flajolet and R. Sedgewick}, \emph{Analytic Combinatorics} (Cambridge University Press, 2009) \cite{FlajoletSedgewick2009}, Part~B --- Chapters~IV--VI develop the complex-analytic foundations used throughout this book: analytic continuation, the transfer theorem, and singularity analysis in full generality. Reading Part~B alongside the present chapter is strongly recommended.

  \item \textbf{P. Flajolet and A.\,M. Odlyzko}, ``Singularity analysis of generating functions,'' \emph{SIAM J.\ Discrete Math.} (1990) \cite{FlajoletOdlyzko1990} --- The original paper introducing the transfer theorem and $\Delta$-domains; more concise than the textbook treatment and useful for seeing the theorem in its original form.

  \item \textbf{L.\,V. Ahlfors}, \emph{Complex Analysis} (3rd ed., McGraw-Hill, 1979) --- The standard graduate-level reference for classical complex analysis. Chapters~1--4 cover all the material in the present chapter (convergence, analyticity, Cauchy's theorem, singularity classification) with full proofs; Chapter~5 introduces residues and contour integration.

  \item \textbf{E.\,C. Titchmarsh}, \emph{The Theory of Functions} (2nd ed., Oxford, 1939) --- A more classical treatment, particularly strong on power series (Chapter~1) and contour integration and residues (Chapter~3); useful as a second reference for the Cauchy--Hadamard theorem and the classification of isolated singularities.

  \item \textbf{P. Henrici}, \emph{Applied and Computational Complex Analysis}, Vol.~1 (Wiley, 1974) --- Covers power series and analytic functions with an emphasis on effective computation and coefficient extraction; bridges the gap between pure complex analysis and the algorithmic side of generating functions.
\end{itemize}
